\section{Introducción}
\label{sec:introduccion}

\subsection{Contexto y Motivación}
La evolución contemporánea de los sistemas de Recuperación Aumentada por Generación (RAG, por sus siglas en inglés) ha transitado desde arquitecturas puramente textuales hacia paradigmas multimodales capaces de procesar, alinear e interpretar información visual, auditiva y textual dentro de un único espacio semántico continuo \cite{prajwal2025multimodal}. En el ámbito de la recuperación de eventos de interés sobre secuencias de vídeo de larga duración (disciplina denominada en la literatura académica como \textit{Long-form Video Understanding} o \textit{Video RAG}), tanto la investigación científica como el desarrollo industrial han experimentado un avance significativo entre los años 2024 y 2026 \cite{multi_rag2025}. Este progreso responde a la necesidad de identificar, localizar y extraer eventos específicos dentro de archivos de vídeo extensos (tales como la mención de un producto, la aparición de un objeto o un cambio de escena), eliminando la necesidad del visionado manual.

\subsection{Planteamiento del Problema}
Históricamente, la arquitectura tradicional para recuperar información en vídeo obligaba a procesar los fotogramas por separado, emplear modelos de transcripción fonética básica para el audio (como Whisper) y extraer texto mediante Reconocimiento Óptico de Caracteres (OCR) para indexarlo todo como texto plano \cite{video_rag2024}. El estado del arte ha demostrado empíricamente que convertir toda la información visual a texto plano genera un cuello de botella informacional y una pérdida severa de contexto temporal, espacial, estructural y emocional, haciendo desaparecer la disposición espacial (\textit{layout}) y las relaciones semánticas entre los objetos visuales \cite{prajwal2025multimodal}. Frente a estas limitaciones, la aparición de modelos de \textit{embeddings} omnimodales permite ingerir texto, vídeo, audio y documentos estructurados (PDF) de forma nativa para proyectarlos en un único vector numérico sin requerir metadatos anotados manualmente \cite{gemini_emb2_docs}.

\subsection{Propuesta y Objetivos del Trabajo Fin de Máster}
El presente Trabajo de Fin de Máster aborda el ciclo de vida completo para el diseño, desarrollo y validación técnica de un sistema avanzado de recuperación de información multimedia capaz de transformar grandes volúmenes de datos no estructurados (vídeo, audio y texto) en representaciones vectoriales densas (\textit{embeddings}) dentro de un espacio semántico unificado. El sistema permite realizar búsquedas semánticas cruzadas (localizando el instante exacto de un evento en un archivo de vídeo a partir de una consulta en lenguaje natural) sin depender de metadatos previos.

Para estructurar este desarrollo, se define el siguiente \textbf{Objetivo General}: diseñar, implementar y validar empíricamente un sistema avanzado de recuperación de información multimedia basado en arquitecturas de vectorización multimodal, capaz de localizar eventos de interés en grandes volúmenes de datos no estructurados a partir de búsquedas semánticas en lenguaje natural. Este propósito se desglosa en cinco \textbf{Objetivos Específicos}:
\newpage
\begin{enumerate}
    \item \textbf{Analizar el estado del arte} tecnológico en el ámbito de los modelos fundacionales multimodales (generación de \textit{embeddings}) y los sistemas de búsqueda semántica (RAG y bases de datos vectoriales), evaluando las alternativas óptimas para el proyecto.
    \item \textbf{Diseñar la arquitectura técnica} y el flujo de procesamiento (\textit{pipeline}) de datos, definiendo los métodos para la ingestión, fragmentación, preprocesamiento y vectorización de los archivos de vídeo, audio y texto.
    \item \textbf{Desarrollar una Prueba de Concepto (PoC)} funcional que integre el modelo de \textit{embeddings} seleccionado, mapeando los distintos tipos de datos en un espacio semántico unificado para permitir búsquedas cruzadas eficientes sin depender de metadatos previos.
    \item \textbf{Evaluar y validar el rendimiento} del sistema mediante la ejecución de un plan de pruebas empíricas, midiendo métricas clave como la precisión en la recuperación de eventos, el coste computacional y la viabilidad técnica general de la solución.
    \item \textbf{Documentar el ciclo de vida completo} del proyecto de ingeniería del software y datos, detallando el marco metodológico utilizado desde la definición de requisitos hasta la validación final.
\end{enumerate}